\documentclass[11pt,a4paper]{article}
% preamble.tex — shared setup for human-verification documents
% Compile each document with:  tectonic <File>.tex

\usepackage[T1]{fontenc}
\usepackage[utf8]{inputenc}
\usepackage{lmodern}
\usepackage{amssymb,amsmath,amsthm}
\usepackage{mathtools}
\usepackage[dvipsnames]{xcolor}
\usepackage{listings}
\usepackage{hyperref}
\usepackage[margin=2.5cm]{geometry}
\usepackage{parskip}

% ── Colours ──────────────────────────────────────────────────────────
\definecolor{leanBg}{HTML}{F7F7F0}
\definecolor{haskellBg}{HTML}{F0F4F7}
\definecolor{keyword}{HTML}{7B3294}
\definecolor{comment}{HTML}{5A8C32}
\definecolor{noteFg}{HTML}{B35900}

% ── Lean 4 listings language ─────────────────────────────────────────
% NOTE: Lean listing blocks use ASCII approximations of unicode symbols
% (e.g. -> for arrow, o. for compose).  The literate key substitutes
% these with proper math-mode renderings in the PDF.
\lstdefinelanguage{Lean4}{
  morekeywords={def,theorem,structure,where,namespace,end,import,fun,
                let,have,show,by,exact,intro,cases,rfl,subst,obtain,
                funext,simp,rw,apply,variable,Type,Prop,instance,
                section,open,class,extends,deriving,sorry,noncomputable},
  sensitive=true,
  morecomment=[l]{--},
  morecomment=[n]{/-}{-/},
  morestring=[b]",
}

% ── Lean listing style ───────────────────────────────────────────────
\lstdefinestyle{lean}{
  language=Lean4,
  backgroundcolor=\color{leanBg},
  basicstyle=\small\ttfamily,
  keywordstyle=\color{keyword}\bfseries,
  commentstyle=\color{comment}\itshape,
  breaklines=true,
  frame=single,
  framerule=0.4pt,
  xleftmargin=4pt,
  xrightmargin=4pt,
  aboveskip=6pt,
  belowskip=4pt,
  literate=
    {->}{{$\to$}}2
    {=>}{{$\Rightarrow$}}2
    {<|}{{$\langle$}}1
    {|>}{{$\rangle$}}1
    {==>}{{$\Longrightarrow$}}3
    { o. }{{$\;\circ\;$}}3
    {forall}{{$\forall$}}1
    {:=}{{$\,:=\,$}}2
    {!=}{{$\neq$}}2
    {<=}{{$\leq$}}2
    {>=}{{$\geq$}}2,
}

% ── Haskell listing style ───────────────────────────────────────────
\lstdefinestyle{haskell}{
  language=Haskell,
  backgroundcolor=\color{haskellBg},
  basicstyle=\small\ttfamily,
  keywordstyle=\color{keyword}\bfseries,
  commentstyle=\color{comment}\itshape,
  breaklines=true,
  frame=single,
  framerule=0.4pt,
  xleftmargin=4pt,
  xrightmargin=4pt,
  aboveskip=6pt,
  belowskip=4pt,
  literate=
    {->}{{$\to$}}2
    {=>}{{$\Rightarrow$}}2,
}

% ── Theorem environments ────────────────────────────────────────────
\theoremstyle{definition}
\newtheorem{defn}{Definition}[section]
\newtheorem{thm}[defn]{Theorem}
\newtheorem{prop}[defn]{Proposition}

\theoremstyle{remark}
\newtheorem*{fidelitynote}{Fidelity note}

% ── Convenience macros ───────────────────────────────────────────────
\newcommand{\lean}[1]{\texttt{#1}}      % inline Lean identifier
\newcommand{\hask}[1]{\texttt{#1}}      % inline Haskell identifier
\newcommand{\Ext}{\mathrm{Ext}}
\newcommand{\Pos}{\mathrm{Pos}}
\newcommand{\Shape}{\mathrm{Shape}}


\title{Human Verification: \texttt{Trajectory}\\[4pt]
  \large Trajectory Systems and Circular Verification}
\author{Machine-generated from \texttt{lean/Trajectory.lean} for review by Neil Ghani}
\date{July 2026}

\begin{document}
\maketitle

\begin{abstract}
This document presents every definition and theorem \emph{statement} from the
Lean~4 module \lean{Trajectory} in three parallel forms: the original Lean,
a readable Haskell rendering, and precise mathematical English.  Proofs are
omitted --- they are machine-checked by Lean.  The module instantiates the
general trust-system theory of \lean{Trust.lean} to \emph{trajectory systems},
where claims are agent events and checking steps are causal dependencies.
The main result is that circular verification collapses canonical trust to
$\bot$ at every event.
\end{abstract}

\tableofcontents
\bigskip

%===================================================================
\section{Overview}
%===================================================================

A \emph{trajectory system} is a trust system (in the sense of
\lean{Containers.Trust}) whose claims are agent actions (\lean{Evt}) and whose
checking steps record causal dependencies between events.  The module
demonstrates two contrasting instances: a \emph{circular} trajectory whose
verification graph has a cycle (MacBeth Day~96) and an \emph{honest} trajectory
with acyclic dependencies.  The canonical trust $\tau$ collapses to
$\bot = \mathtt{unsupported}$ on the circular instance and reaches
\lean{independent} or above on the honest instance.

%===================================================================
\section{The five-level trajectory chain}
%===================================================================

\subsection{Inductive type: \lean{TLevel}}

\begin{defn}[Trajectory trust levels]
The type of trajectory trust strengths is a five-element linearly ordered set:
\[
  \mathtt{unsupported} < \mathtt{shallow} < \mathtt{independent}
  < \mathtt{crossVerified} < \mathtt{mechChecked}.
\]
\end{defn}

\paragraph{Lean.}
\begin{lstlisting}[style=lean]
inductive TLevel where
  | unsupported
  | shallow
  | independent
  | crossVerified
  | mechChecked
deriving DecidableEq, Repr
\end{lstlisting}

\paragraph{Haskell.}
\begin{lstlisting}[style=haskell]
data TLevel
  = Unsupported
  | Shallow
  | Independent
  | CrossVerified
  | MechChecked
  deriving (Eq, Ord, Show)
\end{lstlisting}


\subsection{Definition: \lean{TLevel.le}}

\begin{defn}[Order on \lean{TLevel}]
The order is defined case-by-case: $\mathtt{unsupported} \le b$ for all $b$;
$\mathtt{mechChecked} \le b$ only when $b = \mathtt{mechChecked}$; and so on
through the chain.  The \lean{Chain TLevel} instance packages this into a
decidable linear order with $\bot = \mathtt{unsupported}$.
\end{defn}

\paragraph{Lean.}
\begin{lstlisting}[style=lean]
def TLevel.le : TLevel -> TLevel -> Prop
  | .unsupported, _ => True
  | .shallow, .unsupported => False
  | .shallow, _ => True
  | .independent, .unsupported => False
  | .independent, .shallow => False
  | .independent, _ => True
  | .crossVerified, .mechChecked => True
  | .crossVerified, .crossVerified => True
  | .crossVerified, _ => False
  | .mechChecked, .mechChecked => True
  | .mechChecked, _ => False

instance : Chain TLevel where
  bot := .unsupported
  -- le_refl, le_trans, le_antisymm, le_total, bot_le, dec all proved
\end{lstlisting}

\paragraph{Haskell.}
\begin{lstlisting}[style=haskell]
-- Derived Ord gives the same linear order.
-- bot = Unsupported.
\end{lstlisting}

%===================================================================
\section{Concrete events}
%===================================================================

\subsection{Inductive type: \lean{Evt}}

\begin{defn}[Events]
The five events $e_0, e_1, e_2, e_3, e_4$ are the claims of both the circular
and honest trajectory systems.
\end{defn}

\paragraph{Lean.}
\begin{lstlisting}[style=lean]
inductive Evt where
  | e0 | e1 | e2 | e3 | e4
deriving DecidableEq, Repr
\end{lstlisting}

\paragraph{Haskell.}
\begin{lstlisting}[style=haskell]
data Evt = E0 | E1 | E2 | E3 | E4
  deriving (Eq, Show)
\end{lstlisting}

%===================================================================
\section{The circular trajectory (MacBeth Day 96)}
%===================================================================

\subsection{Definition: \lean{circularTraj}}

\begin{defn}[Circular trajectory system]
A trust system \lean{System Evt TLevel} where each event is its own unique
checking step, and the dependency graph is the chain
$e_0 \leftarrow e_1 \leftarrow e_2 \leftarrow \{e_3, e_4\}$.
The assigned check grades are:
\[
  e_0 \mapsto \mathtt{unsupported},\quad
  e_1, e_2 \mapsto \mathtt{independent},\quad
  e_3 \mapsto \mathtt{shallow},\quad
  e_4 \mapsto \mathtt{unsupported}.
\]
The name ``circular'' refers to the fact that the grades record verification
steps that were logically circular: the supporting evidence ultimately relies
on the claim being supported.
\end{defn}

\paragraph{Lean.}
\begin{lstlisting}[style=lean]
def circularTraj : System Evt TLevel where
  Step := Evt
  concl := id
  prems
    | .e0 => []
    | .e1 => [.e0]
    | .e2 => [.e1]
    | .e3 => [.e2]
    | .e4 => [.e2]
  grade
    | .e0 => .unsupported
    | .e1 => .independent
    | .e2 => .independent
    | .e3 => .shallow
    | .e4 => .unsupported
  stepsInto x := [x]
  mem_stepsInto := ...
  stepsInto_mem f := .head []
\end{lstlisting}

\paragraph{Haskell.}
\begin{lstlisting}[style=haskell]
-- A System Evt TLevel where Step = Evt,
-- prems e0=[], e1=[e0], e2=[e1], e3=[e2], e4=[e2],
-- grade e0=Unsupported, e1=Independent, e2=Independent,
--       e3=Shallow, e4=Unsupported,
-- stepsInto x = [x].
\end{lstlisting}


\subsection{Theorem: \lean{circularTraj\_wf}}

\begin{thm}[Well-foundedness of the circular premise relation]
The premise relation of \lean{circularTraj} is well-founded: the dependency
graph $e_0 \leftarrow e_1 \leftarrow e_2 \leftarrow \{e_3, e_4\}$ is acyclic.
\end{thm}

\paragraph{Lean.}
\begin{lstlisting}[style=lean]
theorem circularTraj_wf : WellFounded circularTraj.Prem
\end{lstlisting}

\paragraph{Haskell.}
\begin{lstlisting}[style=haskell]
-- The dependency DAG is acyclic (e0 has no premises),
-- so the premise relation is well-founded.
\end{lstlisting}

\begin{fidelitynote}
Despite the system being called ``circular,'' the \emph{dependency graph} is
acyclic (the circularity is in the \emph{verification logic}, not the premise
structure).  Well-foundedness is needed to apply the certificate lemma.
\end{fidelitynote}


\subsection{Theorem: \lean{circularTraj\_tau\_bot}}

\begin{thm}[Circular verification collapses trust]
Canonical trust is $\bot$ at every event in the circular trajectory:
\[
  \forall\, e : \mathtt{Evt},\quad
  \tau(\mathtt{circularTraj},\, e) = \mathtt{unsupported}.
\]
The grade degradation propagates from $e_0$ (graded \lean{unsupported}, no
premises) through the chain: each $\tau(e_i)$ is limited by the weakest link
$\tau(e_0) = \bot$.
\end{thm}

\paragraph{Lean.}
\begin{lstlisting}[style=lean]
theorem circularTraj_tau_bot : forall (e : Evt),
    tau circularTraj circularTraj_wf e = .unsupported
\end{lstlisting}

\paragraph{Haskell.}
\begin{lstlisting}[style=haskell]
-- prop_circularBot :: Evt -> Bool
-- prop_circularBot e = tau circularTraj e == Unsupported
\end{lstlisting}

%===================================================================
\section{The honest trajectory}
%===================================================================

\subsection{Definition: \lean{honestTraj}}

\begin{defn}[Honest trajectory system]
The same event set and dependency structure as \lean{circularTraj}, but with
grades reflecting genuine independent verification:
\[
  e_0, e_1, e_2, e_4 \mapsto \mathtt{independent},\quad
  e_3 \mapsto \mathtt{shallow}.
\]
\end{defn}

\paragraph{Lean.}
\begin{lstlisting}[style=lean]
def honestTraj : System Evt TLevel where
  Step := Evt
  concl := id
  prems
    | .e0 => []
    | .e1 => [.e0]
    | .e2 => [.e1]
    | .e3 => [.e2]
    | .e4 => [.e2]
  grade
    | .e0 => .independent
    | .e1 => .independent
    | .e2 => .independent
    | .e3 => .shallow
    | .e4 => .independent
  stepsInto x := [x]
  mem_stepsInto := ...
  stepsInto_mem f := .head []
\end{lstlisting}

\paragraph{Haskell.}
\begin{lstlisting}[style=haskell]
-- Same structure as circularTraj but grade e0 = Independent.
\end{lstlisting}


\subsection{Theorem: \lean{honestTraj\_wf}}

\begin{thm}[Well-foundedness of the honest premise relation]
The premise relation of \lean{honestTraj} is well-founded.
\end{thm}

\paragraph{Lean.}
\begin{lstlisting}[style=lean]
theorem honestTraj_wf : WellFounded honestTraj.Prem
\end{lstlisting}

\paragraph{Haskell.}
\begin{lstlisting}[style=haskell]
-- Same DAG; acyclic and hence well-founded.
\end{lstlisting}


\subsection{Theorems: \lean{honest\_tau\_e0} through \lean{honest\_tau\_e4}}

\begin{thm}[Canonical trust in the honest trajectory]
The canonical trust values in the honest trajectory are:
\begin{align*}
  \tau(e_0) &= \mathtt{independent}, \\
  \tau(e_1) &= \mathtt{independent}, \\
  \tau(e_2) &= \mathtt{independent}, \\
  \tau(e_3) &= \mathtt{shallow}, \\
  \tau(e_4) &= \mathtt{independent}.
\end{align*}
At each event, $\tau$ is the meet of the grade along its unique dependency
chain from $e_0$.  $e_3$ receives only \lean{shallow} because its step is
graded \lean{shallow}; all other events reach \lean{independent}.
\end{thm}

\paragraph{Lean.}
\begin{lstlisting}[style=lean]
theorem honest_tau_e0 :
    tau honestTraj honestTraj_wf .e0 = .independent

theorem honest_tau_e1 :
    tau honestTraj honestTraj_wf .e1 = .independent

theorem honest_tau_e2 :
    tau honestTraj honestTraj_wf .e2 = .independent

theorem honest_tau_e3 :
    tau honestTraj honestTraj_wf .e3 = .shallow

theorem honest_tau_e4 :
    tau honestTraj honestTraj_wf .e4 = .independent
\end{lstlisting}

\paragraph{Haskell.}
\begin{lstlisting}[style=haskell]
-- prop_honestTau :: Evt -> TLevel -> Bool
-- The expected values match the weakest-link computation through the DAG.
\end{lstlisting}

%===================================================================
\section{Translation conventions}
%===================================================================

\begin{enumerate}
\item \textbf{Proof bodies omitted.}  Lean machine-checks all proofs.  Human
  verification covers only the \emph{statements}.

\item \textbf{Chain instance.}  The \lean{instance : Chain TLevel} bundles
  reflexivity, transitivity, antisymmetry, totality, bot-below-all, and
  decidability.  In Haskell this corresponds to \hask{Eq} + \hask{Ord} with
  the natural derived instance.

\item \textbf{Well-foundedness.}  The certificate lemma (\lean{sound\_le\_tau}
  in \lean{Trust.lean}) requires an explicit \lean{WellFounded} proof.  The
  circular trajectory's DAG is acyclic (the circularity is in the verification
  logic, not the graph), so \lean{WellFounded} holds; the Lean proof constructs
  accessibility witnesses node by node.

\item \textbf{Private helper theorems.}  The leaf-to-root computation of
  $\tau$ for each event (\lean{tau\_circ\_e0}, \lean{tau\_circ\_e1}, etc.)
  are \lean{private} helpers leading to the exported
  \lean{circularTraj\_tau\_bot}.  They are stated but not expanded here.
\end{enumerate}

\end{document}
